\begin{summarybox}
\textbf{\textcolor{CSummary}{一句话核心}}

认清函数类型，再检查可导范围；基本公式与求导法则组合即可计算常见函数的导数。

\medskip
\textbf{\textcolor{CSummary}{核心速记表}}
\[
\begin{aligned}
(c)'&=0,\ (x)'=1,\ (x^\alpha)'=\alpha x^{\alpha-1},\\
(1/x)'&=-1/x^2,\ (\sqrt{x})'=1/(2\sqrt{x}),\\
(a^x)'&=a^x\ln a,\ (e^x)'=e^x,\\
(\log_a x)'&=1/(x\ln a),\ (\ln x)'=1/x,\\
(\sin x)'&=\cos x,\ (\cos x)'=-\sin x.
\end{aligned}
\]
\textbf{\textcolor{CSummary}{常见补充（按条件使用）}}
\[
\begin{aligned}
(\sqrt[3]x)'&=1/(3\sqrt[3]{x^2}),\\
(\tan x)'&=1/\cos^2x,\ (\cot x)'=-1/\sin^2x,\\
(\sec x)'&=\sec x\tan x,\\
(\csc x)'&=-\csc x\cot x,\\
(\ln|x|)'&=1/x,\ (|x|)'=x/|x|\ (x\ne0),\\
(x^x)'&=x^x(\ln x+1)\ (x>0).
\end{aligned}
\]
反三角函数见正文拓展，不列入高考核心必记表。

\medskip
\begin{minipage}{\linewidth}
\textbf{\textcolor{CSummary}{三步自检}}
\begin{itemize}[leftmargin=*,itemsep=0.5em,topsep=0.4em]
\item 先看范围：对数真数、根式、分母、幂函数零点；三角函数使用弧度制。
\item 再看结构：和差逐项求导，乘除用法则，复合乘内层导数，变指数先转成指数函数。
\item 最后看符号：余弦、余切、余割和倒数的负号；绝对值折点比较左右差商。
\end{itemize}
\end{minipage}
\end{summarybox}
